% !TeX root = ../thuthesis-example.tex

\chapter{形式化问题建模与理论基础}

为确保本系统设计与算法体系的严密性与科学性，本章对企业级 RAG 知识检索、多模态图谱推理、认知记忆衰减以及工作流状态调度展开数学形式化定义与理论推导。

\section{非结构化知识切分与边界形式化建模}

\begin{definition}[非结构化文档空间]
设自然语言词表为 $\Sigma$。一个非结构化文档 $D$ 形式化表示为有限字符序序列：
\begin{equation}
    D = (c_1, c_2, \ldots, c_N), \quad c_i \in \Sigma, \quad N \in \mathbb{N}^+
\end{equation}
文档集合记为 $\mathcal{D} = \{D_1, D_2, \ldots, D_{|\mathcal{D}|}\}$。
\end{definition}

\begin{definition}[语义切块算子与信息边界约束]
切块算子 $\Phi: \mathcal{D} \rightarrow \mathcal{P}(\mathcal{S})$ 是一个将原始文档映射为一系列语义连续分块集合 $\mathcal{S}_D = \{s_1, s_2, \ldots, s_m\}$ 的分划函数。每一个分块 $s_j$ 由一个三元组严格定义：
\begin{equation}
    s_j = \langle \tau_j, \mathbf{x}_j, \mu_j \rangle
\end{equation}
其中：
\begin{itemize}
    \item $\tau_j \in \Sigma^*$ 为该分块承载的有效文本载荷（Payload）；
    \item $\mathbf{x}_j = [l_j, r_j] \subset [1, N]$ 为分块在原始文档序列中的字符闭区间索引，满足覆盖完备性 $\bigcup_{j=1}^m \mathbf{x}_j = [1, N]$；
    \item $\mu_j = \langle \text{docId}, \text{level}, \text{parent}, \text{breadcrumb} \rangle$ 为分块的层级拓扑元数据。
\end{itemize}
\end{definition}

在实际工程中，分块长度需受严格的物理预算区间约束：
\begin{equation}
    \forall s_j \in \mathcal{S}_D, \quad L_{\min} \le |\tau_j| \le L_{\max}
\end{equation}
设相邻分块间的物理重叠长度为 $\delta = \mathbf{x}_j \cap \mathbf{x}_{j+1}$。重叠机制保证了跨边界实体共现的连续性，满足：
\begin{equation}
    0 \le \delta \le \delta_{\max} < L_{\min}
\end{equation}

\begin{definition}[分块语义内聚度与边界损失]
设句子切分函数为 $\text{Sent}(s) = (u_1, u_2, \ldots, u_k)$，句子向量编码为 $\mathbf{e}(u) \in \mathbb{R}^d$。定义分块 $s_j$ 内部的平均语义内聚度（Intra-Chunk Semantic Cohesion, $\text{SC}$）为相邻句子对余弦相似度的均值：
\begin{equation}
    \text{SC}(s_j) = \frac{1}{k-1} \sum_{i=1}^{k-1} \frac{\mathbf{e}(u_i) \cdot \mathbf{e}(u_{i+1})}{\|\mathbf{e}(u_i)\|_2 \|\mathbf{e}(u_{i+1})\|_2}
\end{equation}
切块算法的全局优化目标是在满足长度约束的前提下，最大化全局语义内聚度期望，并最小化切割导致的上下文断裂惩罚 $\mathcal{L}_{\text{boundary}}$：
\begin{equation}
    \max_{\Phi} \sum_{D \in \mathcal{D}} \left( \sum_{s \in \Phi(D)} \text{SC}(s) - \lambda \mathcal{L}_{\text{boundary}}(\Phi(D)) \right)
\end{equation}
其中 $\lambda > 0$ 为正则化拉格朗日乘子。

\section{多源异构混合检索流形映射模型}

在知识检索阶段，用户查询 $q \in \Sigma^*$ 被同时映射到多个不同维度的数学表示空间中。

\begin{definition}[统一混合表示空间]
定义系统的四路混合检索特征流形元组为 $\mathcal{M} = \langle \mathcal{M}_V, \mathcal{V}_K, \Omega_M, \mathcal{G}_R \rangle$：
\begin{enumerate}
    \item \textbf{稠密语义超球面流形 $\mathcal{M}_V$}：由阿里百炼千问向量编码器 $f_V: \Sigma^* \rightarrow \mathbb{S}^{d-1}$ 生成，$d=1536$，满足 $\|\mathbf{v}\|_2 = 1$。查询向量记为 $\mathbf{q}_V = f_V(q)$，候选切片向量记为 $\mathbf{v}_d = f_V(\tau_d)$。度量方式为超球面余弦测度：
    \begin{equation}
        S_V(q, d) = \langle \mathbf{q}_V, \mathbf{v}_d \rangle = \sum_{i=1}^{d} q_{V,i} \cdot v_{d,i} \in [-1, 1]
    \end{equation}
    
    \item \textbf{稀疏词项倒排空间 $\mathcal{V}_K$}：词项离散空间 $\mathcal{T} = \{t_1, t_2, \ldots, t_{|\mathcal{T}|}\}$。正文由 Tantivy BM25 与 PostgreSQL 三元组（Trigram）及全文检索评分联合表征：
    \begin{equation}
        S_K(q, d) = \alpha \cdot \text{BM25}(q, d) + (1-\alpha) \cdot \left( \max(\text{trgm}(q, d), \text{ts\_rank}(q, d)) \times 0.6 + \text{KW}(q, d) \times 0.4 \right)
    \end{equation}
    其中 $\alpha \in [0, 1]$ 为引擎可用性选择因子。
    
    \item \textbf{结构化属性滤网 $\Omega_M$}：基于一阶谓词逻辑的多租户与时序空间判定：
    \begin{equation}
        \mathbb{I}_M(q, d) = \mathbb{I}(\text{workspace}(d) = w_q) \land \mathbb{I}(\text{day\_no}(d) \in \text{Range}(q)) \in \{0, 1\}
    \end{equation}
    
    \item \textbf{知识拓扑图谱 $\mathcal{G}_R = (\mathcal{V}_E, \mathcal{E}_R, \mathcal{W}_R)$}：其中 $\mathcal{V}_E$ 为实体节点集合，$\mathcal{E}_R \subseteq \mathcal{V}_E \times \mathcal{V}_E$ 为关联边，$\mathcal{W}_R: \mathcal{E}_R \rightarrow \mathbb{R}^+$ 为关系亲密度权重。每个实体关联其起源切片倒排索引 $\text{Segments}(v) \subset \mathcal{S}$。
\end{enumerate}
\end{definition}

\section{倒数排名融合（RRF）单调保序性与弱路径剪枝定理}

在多路混合召回中，各检索路径输出的原始分值度量尺度严重不一致（向量余弦分在 $[0, 1]$ 之间，BM25 分值在 $[0, +\infty)$ 之间，元数据评分在 $[0, 10]$ 之间，图谱评分在 $[0, 12]$ 之间）。

\begin{definition}[倒数排名融合评分函数]
设参与检索的路径集合为 $\mathcal{P} = \{\text{Vector}, \text{Keyword}, \text{Metadata}, \text{Graph}\}$。对于任意有效路径 $p \in \mathcal{P}$，其返回的有序候选序列为 $L_p = (d_{p,1}, d_{p,2}, \ldots, d_{p,|L_p|})$。文档 $d$ 在路径 $p$ 中的排名定义为：
\begin{equation}
    r_p(d) = \begin{cases}
        i, & \text{若 } d = d_{p,i} \\
        +\infty, & \text{若 } d \notin L_p
    \end{cases}
\end{equation}
系统采用带平滑常数 $k \in \mathbb{N}^+$（工程基线固定为 $k=60$）的倒数排名融合（Reciprocal Rank Fusion, RRF）聚合得分：
\begin{equation}
    \text{RRF}(d) = \sum_{p \in \mathcal{P}} \frac{1}{k + r_p(d)}
\end{equation}
\end{definition}

\begin{theorem}[RRF 融合算子的单调保序性与无量纲不变量]
\label{thm:rrf_monotone}
对于任意两个候选文档 $d_A, d_B \in \mathcal{D}$，若在所有参与融合的路径 $p \in \mathcal{P}$ 上均满足 $r_p(d_A) \le r_p(d_B)$，且存在至少一条路径 $p^* \in \mathcal{P}$ 满足严格小于 $r_{p^*}(d_A) < r_{p^*}(d_B)$，则必有：
\begin{equation}
    \text{RRF}(d_A) > \text{RRF}(d_B)
\end{equation}
且对于任意路径分值的严格单调递增变换 $\phi_p: \mathbb{R} \rightarrow \mathbb{R}$，只要其保持局部排名序不变，最终的 RRF 聚合全序关系保持绝对不变。
\end{theorem}

\begin{proof}
由于平滑参数 $k = 60 > 0$，函数 $g(x) = \frac{1}{k + x}$ 在定义域 $[1, +\infty)$ 上严格单调递减，其一阶导数满足：
\begin{equation}
    g'(x) = -\frac{1}{(k + x)^2} < 0, \quad \forall x \ge 1
\end{equation}
由条件已知，$\forall p \in \mathcal{P}, r_p(d_A) \le r_p(d_B)$。利用 $g(x)$ 的严格单调递减性，有：
\begin{equation}
    \frac{1}{k + r_p(d_A)} \ge \frac{1}{k + r_p(d_B)}, \quad \forall p \in \mathcal{P}
\end{equation}
又因为存在 $p^* \in \mathcal{P}$ 使得 $r_{p^*}(d_A) < r_{p^*}(d_B)$，则有严格不等式成立：
\begin{equation}
    \frac{1}{k + r_{p^*}(d_A)} > \frac{1}{k + r_{p^*}(d_B)}
\end{equation}
对所有路径求和，由不等式的可加性直接推得：
\begin{equation}
    \text{RRF}(d_A) = \sum_{p \in \mathcal{P}} \frac{1}{k + r_p(d_A)} > \sum_{p \in \mathcal{P}} \frac{1}{k + r_p(d_B)} = \text{RRF}(d_B)
\end{equation}
此外，由于排名函数 $r_p(d)$ 仅依赖于局部序列的相对比较偏序关系，对原始得分进行任意保持偏序的一致单调变换 $\phi_p$ 均不改变 $r_p(d)$ 的取值，因此 $\text{RRF}(d)$ 具有严格的无量纲尺度不变量特性。
\end{proof}

\begin{lemma}[动态弱路径剪枝判决下界]
\label{lem:weak_path}
设路径 $p$ 的 Top-1 候选在经过全局尺度归一化函数 $\mathcal{N}_p(\cdot) \in [0, 1]$ 处理后的分值为 $\hat{s}_{p,1}$。设弱路径硬阈值为 $\theta_w \in (0, 1)$。若满足：
\begin{equation}
    \hat{s}_{p,1} = \mathcal{N}_p\left(\max_{d \in L_p} S_p(q, d)\right) < \theta_w
\end{equation}
则路径 $p$ 对全局 Top-1 真实候选的置信贡献方差低于白噪声基准。将该路径从候选池排除，使融合退化为在子集 $\mathcal{P} \setminus \{p\}$ 上的投影，可降低混合检索的假阳性率（False Positive Rate）。
\end{lemma}

\section{动态自适应艾宾浩斯记忆强化衰减常微分方程推导}

在智能体的长效认知记忆机制中，信息的留存与遗忘呈现非线性的时序演化特征。

\begin{definition}[动态记忆状态变量]
设一条原子记忆单元为 $m$。在时刻 $t \ge 0$，该记忆的内在巩固强度（Strength）记为 $S(t) \in \mathbb{R}^+$，其瞬时可检索留存率（Retrievability）记为 $R(t) \in [0, 1]$。
\end{definition}

\begin{theorem}[自适应艾宾浩斯强化衰减动力学模型]
\label{thm:ebbinghaus_ode}
设记忆单元在前序历史上被成功检索唤醒的总次数为 $k \in \mathbb{N}$，其固有的重要性先验权重为 $I \in [0, 1]$。记忆留存率 $R(t)$ 服从广义负指数衰减律：
\begin{equation}
    R(t) = \exp\left( -\frac{\Delta t}{S_k} \right)
\end{equation}
其中 $\Delta t = t - t_{\text{last}}$ 为自上一次被唤醒激活以来的时间流逝流。记忆强度 $S_k$ 随唤醒次数 $k$ 的动态演化满足离散强化递推关系：
\begin{equation}
    S_{k+1} = S_k \cdot \left( 1 + \alpha \ln(1 + k) \right), \quad \alpha = 0.20
\end{equation}
其初始强度由重要度调制：$S_0(I) = S_{\text{base}} \cdot (1 + \beta I)$。
\end{theorem}

\begin{proof}
设在连续时间近似下，记忆强度的增量率关于唤醒经历 $k$ 满足饱和增长定律。根据心理物理学 Weber-Fechner 对数感知原理，新刺激对认知记忆印痕的边际强化效应与已发生强化历史的丰富度呈反比衰减：
\begin{equation}
    \frac{\dif S}{\dif k} = \alpha \cdot \frac{S}{1 + k}, \quad \alpha > 0
\end{equation}
对该一阶常微分方程分离变量并积分：
\begin{equation}
    \int \frac{1}{S} \dif S = \alpha \int \frac{1}{1 + k} \dif k \implies \ln S(k) = \alpha \ln(1 + k) + C
\end{equation}
取初始条件 $k=0$ 时 $S(0) = S_0$，可得积分常数 $C = \ln S_0$。两边取指数得到解的连续闭式表达式：
\begin{equation}
    S(k) = S_0 \cdot (1 + k)^\alpha
\end{equation}
在离散工程系统中，考虑单步增量递推，由泰勒一阶展开或差分近似，单步更新乘子满足 $S_{k+1} = S_k \cdot (1 + \alpha \ln(1+k))$。代入留存率方程，半衰期（Half-life, 即 $R(t_{1/2}) = 0.5$ 时所需的时间）为：
\begin{equation}
    t_{1/2}(k) = S_k \cdot \ln 2 = S_0 \cdot (1 + k)^\alpha \cdot \ln 2
\end{equation}
当 $k \rightarrow +\infty$ 时，$t_{1/2}(k) \rightarrow +\infty$，即频繁被检索引用的关键事实具备渐进不遗忘性（Asymptotic Retention）。
\end{proof}

\begin{definition}[多目标 Pareto 综合记忆评分]
综合语义相关性、时序遗忘度、先验重要性以及知识图谱拓扑关联，在时刻 $t$ 的最终长效记忆评分公式定义为多目标凸组合映射：
\begin{equation}
    \text{Score}_{\text{mem}}(m, q, t) = w_1 \cdot \text{Sim}_{\text{cal}}(m, q) + w_2 \cdot R(t) + w_3 \cdot I(m) + w_4 \cdot C_{\text{graph}}(m, q)
\end{equation}
约束权重归一化：$\sum_{i=1}^4 w_i = 1$。其中 $\text{Sim}_{\text{cal}}$ 为经过温度参数 $T=0.15$ 与平衡点 $\tau=0.65$ 校准的 Sigmoid 锐化相似度：
\begin{equation}
    \text{Sim}_{\text{cal}}(m, q) = \frac{1}{1 + \exp\left( -\frac{\cos(f_V(m), f_V(q)) - \tau}{T} \right)}
\end{equation}
该非线性变换大幅拉大了高置信度与中低置信度候选之间的间距，抑制了中等相似度噪声对大模型记忆上下文的干扰。
\end{definition}

\section{基于定界延续（Delimited Continuation）的工作流执行代数模型}

在处理引入人工审批与异步长事务的复杂 DAG 工作流时，系统的状态迁移需满足形式化代数性质。

\begin{definition}[工作流有向无环图]
工作流被形式化为一个带权有向无环图 $\mathcal{G}_{\text{flow}} = (\mathcal{V}_N, \mathcal{E}_D)$，其中 $\mathcal{V}_N$ 为异构节点集合（涵盖任务节点、条件分支网关、大模型推理节点、人工审批挂起节点等），$\mathcal{E}_D \subset \mathcal{V}_N \times \mathcal{V}_N$ 为执行依赖偏序边。
\end{definition}

\begin{definition}[定界延续状态机]
定义工作流运行时的全局状态为元组 $\sigma = \langle \mathcal{E}_X, \mathcal{V}_C, \Gamma, \Theta \rangle$：
\begin{itemize}
    \item $\mathcal{E}_X \in \{\text{READY}, \text{RUNNING}, \text{SUSPENDED}, \text{COMPLETED}, \text{FAILED}\}$ 为全局执行生命周期状态；
    \item $\mathcal{V}_C \subseteq \mathcal{V}_N$ 为当前已收敛完成的节点拓扑闭包；
    \item $\Gamma: \mathcal{V}_N \rightarrow \text{OutputValues}$ 为节点执行结果的键值快照映射；
    \item $\Theta$ 为不可变快照树根指针。
\end{itemize}
定义工作流状态转移函数为 $\tau: \sigma \times \text{Event} \rightarrow \sigma'$。
\end{definition}

\begin{lemma}[非阻塞定界挂起定理]
\label{lem:delimited_continuation}
设节点 $v \in \mathcal{V}_N$ 为人机协同审批节点（HITL Node）。当拓扑调度器遍历至 $v$ 且未检测到预存的审批决议签名时，状态转移函数 $\tau$ 满足：
\begin{equation}
    \tau(\sigma, \text{Exec}(v)) \rightarrow \sigma_{\text{susp}} = \langle \text{SUSPENDED}, \mathcal{V}_C, \Gamma \cup \{v \mapsto \text{Snapshot}\}, \Theta' \rangle
\end{equation}
在此状态跃迁中，调度线程立即退出并调用异步回调返回，无须在内存中维系活跃的执行堆栈或持有着阻塞的物理线程锁。系统将执行上下文冻结为不可变快照树 $\Theta'$ 并持久化至数据库。当外部审批事件 $\text{Approve}(v, \text{Patch})$ 触发时，恢复算子 $\tau_{\text{resume}}$ 通过 CAS 原子操作以 $O(1)$ 时间复杂度恢复至断点继续推进后续子图。
\end{lemma}
